\documentclass[10pt]{article}
\usepackage[utf8]{inputenc}
\usepackage{amsmath,amssymb,amsthm}
\usepackage{geometry}
\geometry{margin=2.5cm}

\newtheorem{theorem}{Theorem}
\newtheorem{lemma}{Lemma}
\newtheorem{definition}{Definition}
\newtheorem{assumption}{Assumption}

\begin{document}

% The following content is intended to replace Sections IV-B and V of the ECPFL manuscript.

\section*{IV-B. Threat Model and Security Goals}

We formalize the security of ECPFL under the standard semi-honest threat model. All cryptographic statements rely on the decisional Ring Learning With Errors assumption and on the pseudorandomness of the public-polynomial generator PRPG.

\subsection*{B.1 System Entities and Trust Assumptions}

ECPFL involves three types of entities.

\begin{enumerate}
    \item \textbf{Key Generation Center (KGC).} The KGC is a fully trusted third party. It generates the public parameters, the per-client encryption keys $\{s_i\}_{i\in[N]}$, the global decryption key $S=\sum_{i=1}^N s_i$, the initial public polynomial $a^{(0)}$, and the PRPG seed $B$. It distributes $s_i$, $S$, $a^{(0)}$, and $B$ to client $i$ through secure channels during system initialization. Because the KGC must generate compensatory ciphertexts on behalf of inactive clients, it is assumed to retain $\{s_i\}_{i\in[N]}$ throughout the protocol execution. The KGC is assumed not to collude with the server or with any subset of clients.

    \item \textbf{Server.} The server is an honest-but-curious central aggregator. It faithfully executes the prescribed aggregation, consensus voting, queue maintenance, and prefetch scheduling procedures, but it may attempt to infer private information from all messages and metadata it observes.

    \item \textbf{Clients.} Each client is an honest-but-curious data owner. Clients execute local training, encryption, decryption, and residual update procedures honestly, but a client may attempt to infer other clients' private gradients from the decrypted global updates and from any protocol messages it receives.
\end{enumerate}

\subsection*{B.2 Adversarial Capabilities}

The adversary may corrupt the server, any strict subset of clients of size at most $N-2$, or external passive eavesdroppers that observe all messages transmitted over public channels. The corrupted parties follow the protocol but pool their views to infer private information. The KGC is not corrupted. We assume secure channels exist only between the KGC and the clients for the one-time distribution of keys and the seed $B$; all other communication occurs over public channels observable by the server and by external eavesdroppers.

The server knows the following information by design: the active set $S(t)$, the inactive set $\bar{S}(t)$, the prefetch set $P(t^*)$, the consensus set $K(t)$, and the plaintext vote masks $\{v_i(t)\}_{i\in S(t)}$. Therefore, the security analysis does not treat these quantities as secrets against the server. External eavesdroppers, however, observe only the ciphertext traffic on public channels.

\subsection*{B.3 Security Goals}

ECPFL aims to achieve the following security guarantees.

\begin{enumerate}
    \item \textbf{Ciphertext confidentiality.} Each uploaded ciphertext is semantically secure under a chosen-plaintext attack. Formally, the encryption primitive is IND-CPA secure.

    \item \textbf{Server obliviousness.} The server cannot learn the plaintext of any individual client's gradient or the plaintext of the aggregated global update, despite observing the ciphertexts, the consensus sets, and the client-selection metadata.

    \item \textbf{Collusion resistance.} Even if up to $N-2$ clients collude with the server and pool their keys and transcripts, they cannot recover the private plaintext of any remaining honest client.

    \item \textbf{Compensatory ciphertext indistinguishability.} A compensatory ciphertext generated by the KGC for an inactive client is computationally indistinguishable from a valid client ciphertext to any external eavesdropper.

    \item \textbf{Cross-round ciphertext privacy.} The prefetch aggregate ciphertext spanning $R$ rounds is computationally indistinguishable from a random element to any external eavesdropper, and an authorized client that decrypts it learns only the sum of the $R$ rounds' updates, not the per-round components.
\end{enumerate}

\section*{V. Security Analysis}

The security of ECPFL is established by the following sequence of theorems. All statements hold under the dRLWE assumption with public parameters $(n,q,\chi_s,\chi_e)$ and under the pseudorandomness of PRPG.

\begin{theorem}[IND-CPA Security]
The encryption scheme $\mathsf{Enc}(m;s_i,a^{(t)})=[a^{(t)}\cdot s_i+\tilde{p}\cdot e+m]_q$ with $e\leftarrow\chi_e$ is IND-CPA secure. For any probabilistic polynomial-time adversary $\mathcal{A}$,
\[
\bigl|\Pr[\mathsf{Exp}^{\mathsf{ind\text{-}cpa}}_{\mathcal{A}}(\lambda)=1]-\tfrac{1}{2}\bigr|\le\mathsf{negl}(\lambda).
\]
\end{theorem}

\begin{proof}
Consider the IND-CPA experiment for a challenge client $i^*$. The adversary receives the public parameters, the seed $B$, and may obtain encryptions under any client key of its choice. It then submits two challenge plaintexts $m_0,m_1\in R_p$ and receives $c^*=[a^{(t)}\cdot s_{i^*}+\tilde{p}\cdot e^*+m_b]_q$ for a random bit $b$.

Define a sequence of hybrid distributions over $(B,a^{(0)},a^{(t)},c^*)$.

\textbf{Hybrid 0.} This is the real experiment. The public polynomial is generated as $a^{(t)}\leftarrow\mathsf{PRPG}(B,t)$.

\textbf{Hybrid 1.} The public polynomial $a^{(t)}$ is replaced by a uniformly random element $u\leftarrow U(R_q)$. By the pseudorandomness of PRPG, Hybrid 0 and Hybrid 1 are computationally indistinguishable.

\textbf{Hybrid 2.} The challenge ciphertext is replaced by $c^*=[u'+m_b]_q$ where $u'\leftarrow U(R_q)$. By the dRLWE assumption, the pair $(a^{(t)},a^{(t)}\cdot s_{i^*}+\tilde{p}\cdot e^*)$ is computationally indistinguishable from $(a^{(t)},u')$. Adding the fixed plaintext $m_b$ preserves this indistinguishability, so Hybrid 1 and Hybrid 2 are computationally indistinguishable.

In Hybrid 2, $c^*$ is uniformly random and independent of $b$, so the adversary's advantage is exactly zero. Therefore, the advantage in the real experiment is negligible.
\end{proof}

\begin{theorem}[Server Obliviousness]
The server cannot compute the plaintext of any individual client's gradient or the plaintext of the aggregated global update from its protocol view.
\end{theorem}

\begin{proof}
The server's view consists of the active client ciphertexts $\{ct_{i,j}^{(t)}\}_{i\in S(t),j\in K(t)}$, the compensatory ciphertexts $\{ct_{k,j}^{(t)}\}_{k\in\bar{S}(t),j\in K(t)}$ received from the KGC, the vote masks $\{v_i(t)\}_{i\in S(t)}$, and the selection sets $S(t)$, $\bar{S}(t)$, $P(t^*)$, $K(t)$.

Each active client ciphertext has the form $ct_{i,j}^{(t)}=[a^{(t)}\cdot s_i+\tilde{p}\cdot e_{i,j}^{(t)}+pt_{i,j}^{(t)}]_q$. Under the dRLWE assumption, the masking term $(a^{(t)},a^{(t)}\cdot s_i+\tilde{p}\cdot e_{i,j}^{(t)})$ is computationally indistinguishable from a uniform random pair. Hence the server learns no information about $pt_{i,j}^{(t)}$ from the ciphertext.

The aggregated ciphertext for pack $j$ is $C_j(t)=[a^{(t)}\cdot S+\tilde{p}\cdot E_j(t)+m_j(t)]_q$, where $S=\sum_{i=1}^N s_i$. Recovering the aggregate plaintext $m_j(t)$ requires knowledge of the global decryption key $S$. The server is not given $S$ and does not possess sufficient key shares to compute it unless it colludes with at least $N-1$ clients, which is excluded by the threat model. Therefore, the server remains oblivious to both individual and aggregated plaintexts.
\end{proof}

\begin{theorem}[Collusion Resistance]
ECPFL is secure against collusion among any $\kappa<N-1$ clients, that is, at most $N-2$ clients, together with the server. Even under such collusion, the plaintexts of the remaining honest clients remain computationally hidden.
\end{theorem}

\begin{proof}
Suppose $\kappa=N-2$ clients collude with the server. Without loss of generality, let clients $1,\dots,N-2$ be corrupted and let clients $N-1$ and $N$ remain honest. The colluders know their own keys $\{s_k\}_{k=1}^{N-2}$, the global key $S$, and the aggregate plaintext $m_j(t)$ decrypted with $S$.

From $S=\sum_{i=1}^N s_i$, the colluders can compute $s_{N-1}+s_N=S-\sum_{k=1}^{N-2}s_k$. They also know the aggregate plaintext $m_j(t)$ because they possess $S$, and they can subtract their own plaintexts to obtain $pt_{N-1,j}^{(t)}+pt_{N,j}^{(t)}$. However, they face the following system of equations:
\begin{align*}
ct_{N-1,j}^{(t)}&=[a^{(t)}\cdot s_{N-1}+\tilde{p}\cdot e_{N-1,j}^{(t)}+pt_{N-1,j}^{(t)}]_q,\\
ct_{N,j}^{(t)}&=[a^{(t)}\cdot s_N+\tilde{p}\cdot e_{N,j}^{(t)}+pt_{N,j}^{(t)}]_q,\\
s_{N-1}+s_N&=S-\sum_{k=1}^{N-2}s_k,\\
pt_{N-1,j}^{(t)}+pt_{N,j}^{(t)}&=m_j(t)-\sum_{k=1}^{N-2}pt_{k,j}^{(t)}.
\end{align*}
This system contains six unknowns $(s_{N-1},s_N,e_{N-1,j}^{(t)},e_{N,j}^{(t)},pt_{N-1,j}^{(t)},pt_{N,j}^{(t)})$ but only four effective constraints. The individual secret key $s_{N-1}$ is not determined by the sum $s_{N-1}+s_N$ alone; under the dRLWE assumption, the pair $(a^{(t)},a^{(t)}\cdot s_{N-1}+\tilde{p}\cdot e_{N-1,j}^{(t)})$ remains computationally indistinguishable from random even when the sum $s_{N-1}+s_N$ is known. Consequently, decrypting $ct_{N-1,j}^{(t)}$ to obtain $pt_{N-1,j}^{(t)}$ is infeasible. The same argument applies to client $N$.

If $\kappa=N-1$ clients collude with the server, they can compute $s_h=S-\sum_{\ell\neq h}s_\ell$ for the single remaining honest client $h$ and decrypt its individual ciphertext. Thus the guarantee holds only for $\kappa<N-1$.
\end{proof}

\begin{theorem}[Compensatory Ciphertext Indistinguishability]
A compensatory ciphertext $ct_{k,j}^{(t)}=[a^{(t)}\cdot s_k+\tilde{p}\cdot e_{k,j}^{(t)}]_q$ generated by the KGC for an inactive client $k$ is computationally indistinguishable from a valid client ciphertext encrypting an arbitrary plaintext $m\in R_p$ to any external eavesdropper observing only public-channel traffic.
\end{theorem}

\begin{proof}
Consider the two distributions
\begin{align*}
\mathcal{D}_0 &= \{(a^{(t)},ct_{k,j}^{(t)}):ct_{k,j}^{(t)}=[a^{(t)}\cdot s_k+\tilde{p}\cdot e_{k,j}^{(t)}]_q\},\\
\mathcal{D}_1 &= \{(a^{(t)},ct_{k,j}^{(t)}):ct_{k,j}^{(t)}=[a^{(t)}\cdot s_k+\tilde{p}\cdot e_{k,j}^{(t)}+m]_q\}.
\end{align*}
By the dRLWE assumption, the base tuple $(a^{(t)},a^{(t)}\cdot s_k+\tilde{p}\cdot e_{k,j}^{(t)})$ is computationally indistinguishable from a uniform random pair. Adding a fixed plaintext $m$ to the second component preserves this indistinguishability. Hence $\mathcal{D}_0$ and $\mathcal{D}_1$ are computationally indistinguishable. The same argument applies when comparing a compensatory ciphertext to the ciphertext of any active client.

Note that this indistinguishability holds against external eavesdroppers. The server itself knows which clients are inactive because it selects $S(t)$ and $\bar{S}(t)$, so the theorem does not claim dropout-status hiding from the server.
\end{proof}

\begin{theorem}[Cross-Round Ciphertext Privacy]
Let $C_j(t+r)=[a^{(t+r)}\cdot S+\tilde{p}\cdot E_j(t+r)+m_j(t+r)]_q$ for $r=0,\dots,R-1$. The prefetch aggregate
\[
C_{\mathsf{pre},j}=\sum_{r=0}^{R-1}C_j(t+r)=\bigl[A_{\mathsf{pre}}\cdot S+\tilde{p}\cdot E_{\mathsf{pre}}+m_{\mathsf{pre},j}\bigr]_q
\]
with $A_{\mathsf{pre}}=\sum_{r=0}^{R-1}a^{(t+r)}$ and $E_{\mathsf{pre}}=\sum_{r=0}^{R-1}E_j(t+r)$ satisfies the following two properties.
\begin{enumerate}
    \item[(i)] The pair $(A_{\mathsf{pre}},C_{\mathsf{pre},j})$ is computationally indistinguishable from $(A_{\mathsf{pre}},u)$ where $u\leftarrow U(R_q)$.
    \item[(ii)] An authorized client that decrypts $C_{\mathsf{pre},j}$ with $S$ obtains only the $R$-round sum $m_{\mathsf{pre},j}=\sum_{r=0}^{R-1}m_j(t+r)$ and learns no per-round component $m_j(t+r)$ beyond what is implied by the sum.
\end{enumerate}
\end{theorem}

\begin{proof}
For (i), each $a^{(t+r)}$ is pseudorandom by the PRPG assumption, so their sum $A_{\mathsf{pre}}$ is pseudorandom. The term $A_{\mathsf{pre}}\cdot S+\tilde{p}\cdot E_{\mathsf{pre}}$ is an RLWE sample with secret $S$ and error $E_{\mathsf{pre}}$, which masks the plaintext $m_{\mathsf{pre},j}$ under the dRLWE assumption. Therefore $(A_{\mathsf{pre}},C_{\mathsf{pre},j})$ is computationally indistinguishable from $(A_{\mathsf{pre}},u)$.

For (ii), decryption yields
\[
m_{\mathsf{pre},j}=\bigl[C_{\mathsf{pre},j}-A_{\mathsf{pre}}\cdot S\bigr]_{q,p}=\sum_{r=0}^{R-1}m_j(t+r).
\]
The client obtains exactly the sum of the $R$ per-round plaintexts. The cryptographic protocol reveals no per-round ciphertext or plaintext to the client; only the $R$-round aggregate is disclosed. Therefore, decomposing $m_{\mathsf{pre},j}$ into its individual components would require information not provided by the protocol.
\end{proof}

\end{document}
